\chapter{PHƯƠNG PHÁP}

\section{Định nghĩa bài toán}
Giả sử có tập dữ liệu
\[
D=\{(x_i,y_i)\},\quad i=1,\ldots,N,
\]
gồm các cặp $(x_i,y_i)$, trong đó $x_i\in\mathbb{R}^d$ là vector các đặc trưng đầu vào của bệnh nhân, $y_i\in\{0,1\}$ là biến nhị phân, với 1 biểu thị có đột quỵ và 0 biểu thị không đột quỵ; $N$ là số bản ghi và $d$ là số đặc trưng đầu vào. Biến mục tiêu $y_i$ được sử dụng để huấn luyện và kiểm thử mô hình học máy.

\section{Tiền xử lý dữ liệu}
Trước khi xử lý dữ liệu, kiểm tra bản ghi trùng được thực hiện nhằm bảo đảm tính duy nhất của tập dữ liệu và tránh làm sai lệch kết quả mô hình do các mẫu lặp. Các giá trị thiếu (missing) chỉ xuất hiện ở cột BMI và được điền bằng giá trị trung bình, giúp duy trì phân bố của cột và tránh mô hình bị thiên lệch do dữ liệu không đầy đủ. Ngoại lai (outliers) được kiểm tra bằng boxplot kết hợp phương pháp khoảng tứ phân vị (IQR): các giá trị nằm ngoài 1.5 lần IQR được xác định là bất thường và thay bằng giá trị biên tương ứng. Những điều chỉnh này được thực hiện để bảo đảm tính nhất quán của dữ liệu và ngăn các bất thường làm sai lệch kết quả mô hình.

\noindent Phân tích dữ liệu khám phá (EDA) [28] được sử dụng để phân tích sơ bộ và xác định các vấn đề tiềm ẩn, chẳng hạn mất cân bằng lớp và dữ liệu thiếu, có thể ảnh hưởng tiêu cực đến mô hình.

\section{PCA song song hóa bảng OpenMP}

Dữ liệu được tâm hóa bằng cách xây dựng ma trận
\[
X_{\text{centered}} = X-\bar{x},
\qquad
\bar{x}=\frac{1}{N}\sum_{i=1}^{N}x_i,
\]
trong đó $\bar{x}$ là vector giá trị trung bình.

\noindent Ma trận hiệp phương sai $S$ được tính như sau:
\[
S=\frac{1}{N-1}X_{\text{centered}}^{T}X_{\text{centered}}.
\]

\noindent Các trị riêng và vector riêng của ma trận hiệp phương sai $S$ được tính bằng phương pháp Jacobi, được song song hóa bằng OpenMP để tăng tốc:
\[
Sv_j=\lambda_jv_j,
\]
trong đó $\lambda_j$ là các trị riêng và $v_j$ là các vector riêng.

\noindent Để giảm chiều, $k$ vector riêng lớn nhất $V_k$ được chọn, sau đó dữ liệu được biến đổi thành ma trận đặc trưng mới:
\[
X_{\text{reduced}}=X_{\text{centered}}V_k.
\]

\section{Mô hình XGBoost}

\noindent Ma trận đặc trưng $X_{\text{reduced}}$ được đưa vào thuật toán XGBoost. Hàm mục tiêu của mô hình có dạng:
\[
L(\theta)=\sum_{i=1}^{N}\ell\left(y_i,\hat y_i(\theta)\right)+\Omega(\theta),
\]
trong đó
\[
\ell(y_i,\hat y_i)=-\left[y_i\log(\hat y_i)+(1-y_i)\log(1-\hat y_i)\right]
\]
là hàm mất mát logistic, $\hat y_i(\theta)$ là dự báo của mô hình, $\theta$ là các tham số mô hình và $\Omega(\theta)$ là thành phần regularization.

\noindent Để tối ưu hiệu năng mô hình, chúng tôi sử dụng Grid Search để điều chỉnh siêu tham số. Quy trình này tìm kiếm có hệ thống trên một tập siêu tham số xác định trước nhằm tối thiểu hóa hàm mất mát và cải thiện độ chính xác. Trong trường hợp này, các siêu tham số của XGBoost được điều chỉnh gồm độ sâu cây, learning rate, số lượng cây, gamma (regularization), minimum child weight, tỷ lệ dữ liệu huấn luyện và tỷ lệ đặc trưng. Mô hình được đánh giá bằng ROC-AUC, phù hợp với bài toán phân loại nhị phân và cho biết khả năng phân biệt giữa các lớp. Mục tiêu là tìm sự cân bằng tối ưu giữa độ chính xác huấn luyện và khả năng khái quát hóa, ngăn overfitting trong khi cải thiện hiệu năng mô hình.

\section{Cân bằng lớp dữ liệu}
\noindent Các phương pháp cân bằng sau được sử dụng:
\begin{itemize}
  \item Undersampling để giảm số lượng mẫu của lớp chiếm ưu thế.
  \item Oversampling bằng SMOTE (Synthetic Minority Over-sampling Technique) [29] để sinh các mẫu tổng hợp cho lớp ít đại diện hơn.
\end{itemize}

\section{Diễn giải mô hình bằng SHAP}
\noindent Để diễn giải dự báo của mô hình, chúng tôi sử dụng phương pháp SHAP [30], ước lượng đóng góp của từng đặc trưng $j$ vào dự báo của từng mẫu $i$:
\[
\operatorname{SHAP}_{j}(x_i)=\phi_j(x_i),
\]
trong đó $\phi_j(x_i)$ là giá trị SHAP thể hiện đóng góp của đặc trưng $j$ vào dự báo cho mẫu $i$. Ở đây $i$ là chỉ số của các mẫu riêng lẻ trong tập dữ liệu, $i=1,2,\ldots,N$, và $N$ là số bản ghi; $j$ là chỉ số của các đặc trưng riêng lẻ, $j=1,2,\ldots,d$, với $d$ là số đặc trưng đầu vào của mỗi bệnh nhân. SHAP được áp dụng cho các đặc trưng gốc chứ không phải các thành phần chính, nhằm duy trì khả năng diễn giải dự báo mô hình theo các biến đầu vào ban đầu.

\noindent Hình~\ref{fig:flow} minh họa quy trình của phương pháp đề xuất. Thuật toán 1 mô tả cách triển khai tìm trị riêng và vector riêng bằng thuật toán quay Jacobi sử dụng OpenMP.

\paperfig[0.58\linewidth]{figures_original/fig-005.jpg}{Sơ đồ luồng của phương pháp đề xuất.\label{fig:flow}}

\noindent\textbf{Thuật toán 1: Tính trị riêng và vector riêng song song}

\noindent Thuật toán trình bày phép quay Jacobi song song. Ma trận đầu vào được khởi tạo cùng ma trận vector riêng, sau đó lặp cho đến khi hội tụ hoặc đạt số vòng lặp tối đa. Ở mỗi vòng lặp, thuật toán tìm phần tử ngoài đường chéo có trị tuyệt đối lớn nhất, tính góc quay Jacobi, cập nhật ma trận và vector riêng, đồng thời sử dụng các chỉ thị \#pragma omp để song song hóa những vòng lặp độc lập. Các vùng tới hạn (critical) được dùng khi nhiều luồng có khả năng cập nhật chung một biến, nhằm ngăn race condition. Thuật toán kiểm tra hội tụ bằng cách theo dõi ngưỡng chính xác eps, cho phép thoát hiệu quả khỏi vòng lặp khi đạt độ chính xác mong muốn. Sau hội tụ, các phần tử đường chéo của ma trận chứa các trị riêng và ma trận tích lũy chứa các vector riêng.

\noindent Việc lựa chọn phương pháp Jacobi xuất phát từ tính phù hợp của nó đối với song song hóa khi làm việc với các tập dữ liệu nhiều chiều. Mặc dù các phương pháp thay thế như phân rã QR hoặc power iteration cũng có thể hiệu quả, phương pháp Jacobi được chọn vì tính đơn giản và khả năng cân bằng giữa độ chính xác và hiệu năng. Trong nghiên cứu tương lai, nhóm tác giả dự định so sánh Jacobi với các phương pháp thay thế này.
Triển khai PCA bằng OpenMP cải thiện đáng kể tốc độ xử lý các tập dữ liệu lớn bằng cách song song hóa các bước then chốt như tính trị riêng và vector riêng của ma trận hiệp phương sai, phân phối tải trên nhiều luồng và giảm đáng kể thời gian thực thi.

\section{Tổng kết phương pháp}

Tóm lại, phương pháp đề xuất xây dựng một pipeline gồm ba khối chính: XGBoost để dự báo, PCA (song song hóa bằng OpenMP) để giảm chiều, và SHAP để diễn giải kết quả.

\noindent XGBoost được lựa chọn vì độ tin cậy trên các tập dữ liệu mất cân bằng, nơi các trường hợpthuộc lớp thiểu số xuất hiện phổ biến như trong bài toán này. PCA giúp giảm chiều dữ liệu ytế nhiều chiều, giữ lại thông tin quan trọng nhất đồng thời loại bỏ dư thừa và nhiễu. SHAP cung cấp cách diễn giải dự báo theo các biến có ý nghĩa lâm sàng, giúp bác sĩ hiểu các yếutố khác nhau tác động thế nào đến nguy cơ đột quỵ. Sự kết hợp này cho phép xử lý các tập dữ liệu lớn, phức tạp, hỗ trợ bác sĩ ra quyết định dựa trên sự giải thích của mô hình, đồng thời giảm nguy cơ overfitting và cải thiện hiệu năng dự báo.